\documentclass[UTF8,12pt,a4paper]{ctexart}
\usepackage[margin=2.2cm]{geometry}
\usepackage{amsmath,amssymb,booktabs,longtable,array,xcolor,hyperref,seqsplit}
\usepackage{enumitem}
\setlist{nosep,leftmargin=2em}
\hypersetup{colorlinks=true,linkcolor=blue!50!black,urlcolor=blue!60!black}
\setlength{\parindent}{2em}
\setlength{\parskip}{0.35em}
\setlength{\emergencystretch}{3em}
\hbadness=10000
\sloppy
\setcounter{tocdepth}{3}
\title{C题建模总思路}
\author{}
\date{}
\begin{document}
\maketitle
\tableofcontents
\clearpage
\begin{center}\LARGE\bfseries C题总思路表\end{center}
\begin{quote}用途：作为论文写作底稿，记录最终采用的数据处理、模型、训练、验证和结论口径。\end{quote}
\section{一、总体路线}
\begingroup
\footnotesize
\setlength{\tabcolsep}{1.5pt}
\begin{longtable}{|>{\raggedright\arraybackslash}p{0.280\textwidth}|>{\raggedright\arraybackslash}p{0.280\textwidth}|>{\raggedright\arraybackslash}p{0.280\textwidth}|}
\hline
阶段 & 核心任务 & 当前状态 \\ \hline
\endfirsthead
\hline
阶段 & 核心任务 & 当前状态 \\ \hline
\endhead
1. 数据清洗 & 修正数据格式、处理重复检测、删失信息、缺失值及异常字段 & 完善中 \\ \hline
2. 数据结构构造 & 分别形成纵向记录数据和患者级首次达标数据 & 已完成 \\ \hline
3. 模型构建 & 关系模型、首次达标时间模型及多因素模型 & 已完成 \\ \hline
4. BMI分组与时点优化 & 搜索分组切点和各组推荐检测孕周 & 已完成 \\ \hline
5. 验证与稳健性 & 嵌套验证、Bootstrap、误差敏感性和泄漏审计 & 已完成 \\ \hline
6. 最终建议 & 汇总结论、适用范围和限制 & 已完成 \\ \hline
\end{longtable}
\endgroup

\section{二、数据清洗}
\begingroup
\scriptsize
\setlength{\tabcolsep}{1.5pt}
\begin{longtable}{|>{\raggedright\arraybackslash}p{0.140\textwidth}|>{\raggedright\arraybackslash}p{0.140\textwidth}|>{\raggedright\arraybackslash}p{0.140\textwidth}|>{\raggedright\arraybackslash}p{0.140\textwidth}|>{\raggedright\arraybackslash}p{0.140\textwidth}|>{\raggedright\arraybackslash}p{0.140\textwidth}|}
\hline
编号 & 数据问题 & 当前处理方案 & 处理对象/阶段 & 处理理由 & 状态 \\ \hline
\endfirsthead
\hline
编号 & 数据问题 & 当前处理方案 & 处理对象/阶段 & 处理理由 & 状态 \\ \hline
\endhead
1 & 同一孕妇存在多次检测，并存在短期内多次检测。男胎1082条检测记录对应267名孕妇；女胎605条检测记录对应147名孕妇。 & 问题1保留全部纵向检测记录，以孕妇随机截距处理同一孕妇记录间的相关性；问题2、3按孕妇转换为一条患者级记录。同一孕妇同一孕周的重复Y浓度取中位数，不同孕周记录按时间排序。训练、验证、测试划分及Bootstrap均以孕妇为单位。 & 重复检测、数据划分与建模单位 & 既保留问题1需要的纵向变化信息，又避免重复计权、伪独立和同一孕妇跨训练集与测试集造成的信息泄漏。 & 已确认 \\ \hline
2 & 原始孕周采用类似 \textsf{\seqsplit{12w+3}} 的字符串格式，不能直接用于数值计算。 & 将 \textsf{\seqsplit{w+d}} 转换为十进制周：\textsf{\seqsplit{w+d/7}}。例如 \textsf{\seqsplit{12w+3=12+3/7=12.4286}} 周。 & 男胎、女胎全部检测记录 & 统一孕周量纲，便于排序、计算时间间隔、构造删失区间及模型训练。 & 已确认 \\ \hline
3 & 检测时间不能精确确定Y染色体浓度首次达到4\%的真实时间，存在左删失、区间删失和右删失。 & 同一孕妇同一孕周的重复Y浓度取中位数并按孕周排序。第一次检测已经达到4\%时构造左删失区间 \textsf{\seqsplit{[0,U]}}；前一次未达标、后一次达标时构造区间删失 \textsf{\seqsplit{[L,U]}}；所有检测均未达标时构造右删失区间 \textsf{\seqsplit{[L,+∞)}}。不删除删失样本，也不把某次检测孕周直接作为精确首次达标时间。 & 男胎问题2、问题3的首次达标时间 & 保留首次达标时间的不确定性，避免把检测时点误当成真实达标时点，并充分利用尚未达标患者的信息。 & 已确认 \\ \hline
4 & 部分字段存在条目缺失。男胎末次月经缺失12条、非整倍体字段空白956条；女胎BMI缺失1条、两列无名列全部为空，非整倍体字段大量空白。 & 沿用现有分类处理：问题1—3不使用的缺失字段保留但不进入模型；关键建模字段不做隐式填补，缺失或无法转换时停止运行；女胎两列全空无名列删除，BMI在每个训练折内用该折中位数填补并增加“BMI缺失”指示变量；女胎非整倍体空白按“未列出T13/T18/T21异常”构造为0标签。 & 男胎、女胎原始数据及各问题建模输入 & 避免为无关字段制造虚假数值；防止利用测试集统计量填补造成泄漏；保留BMI缺失本身可能携带的信息；按附件标签语义区分“空白输出”和普通数值缺失。 & 已确认 \\ \hline
\end{longtable}
\endgroup

\section{已确认规则}
\subsection{规则 DC-01：孕周格式统一}
设原始孕周为 \textsf{\seqsplit{w+d}}，转换后的十进制孕周为

\[
W=w+\frac{d}{7}.
\]
示例：

\begingroup
\footnotesize
\setlength{\tabcolsep}{1.5pt}
\begin{longtable}{|>{\raggedright\arraybackslash}p{0.420\textwidth}|>{\raggedright\arraybackslash}p{0.420\textwidth}|}
\hline
原始孕周 & 十进制孕周 \\ \hline
\endfirsthead
\hline
原始孕周 & 十进制孕周 \\ \hline
\endhead
\textsf{\seqsplit{12w+0}} & 12.0000 \\ \hline
\textsf{\seqsplit{12w+3}} & 12.4286 \\ \hline
\textsf{\seqsplit{15w+6}} & 15.8571 \\ \hline
\end{longtable}
\endgroup

执行要求：无法匹配 \textsf{\seqsplit{w+d}} 形式、缺少孕周或天数超出合理范围时，不进行猜测性转换，应记录并进入异常检查。

\subsection{规则 DC-02：重复检测分层处理}
\begingroup
\footnotesize
\setlength{\tabcolsep}{1.5pt}
\begin{longtable}{|>{\raggedright\arraybackslash}p{0.420\textwidth}|>{\raggedright\arraybackslash}p{0.420\textwidth}|}
\hline
情形 & 处理方式 \\ \hline
\endfirsthead
\hline
情形 & 处理方式 \\ \hline
\endhead
问题1的不同孕周重复检测 & 全部保留，使用孕妇随机截距描述同一孕妇记录间的相关性。 \\ \hline
问题2、3的同孕妇同孕周重复检测 & Y染色体浓度取中位数，避免单次测量波动和重复计权。 \\ \hline
问题2、3的同孕妇不同孕周检测 & 按十进制孕周排序，用于判断未达标—达标的时间区间。 \\ \hline
训练、验证和测试划分 & 按孕妇划分，同一孕妇的记录不得跨集合。 \\ \hline
Bootstrap & 按孕妇抽样，不把同一孕妇的多条记录独立抽样。 \\ \hline
\end{longtable}
\endgroup

处理后的数据单位：问题1仍为1082条纵向记录、267名孕妇；问题2、3为267条患者级记录。

\subsection{规则 DC-03：首次达到4\%的删失区间}
在完成同孕妇同孕周中位数聚合后，按十进制孕周从早到晚排列，以 \textsf{\seqsplit{Y染色体浓度$\geq$0.04}} 为达标标准。

\begingroup
\footnotesize
\setlength{\tabcolsep}{1.5pt}
\begin{longtable}{|>{\raggedright\arraybackslash}p{0.210\textwidth}|>{\raggedright\arraybackslash}p{0.210\textwidth}|>{\raggedright\arraybackslash}p{0.210\textwidth}|>{\raggedright\arraybackslash}p{0.210\textwidth}|}
\hline
删失类型 & 观测情况 & 区间构造 & 当前人数 \\ \hline
\endfirsthead
\hline
删失类型 & 观测情况 & 区间构造 & 当前人数 \\ \hline
\endhead
左删失 & 第一次检测已经达到4\% & \textsf{\seqsplit{T$\in$[0,U]}}，其中 \textsf{\seqsplit{U}} 为第一次检测孕周 & 217 \\ \hline
区间删失 & 前一次未达到4\%，后一次达到4\% & \textsf{\seqsplit{T$\in$[L,U]}}，其中 \textsf{\seqsplit{L}} 为最后一次未达标孕周，\textsf{\seqsplit{U}} 为第一次达标孕周 & 43 \\ \hline
右删失 & 所有检测均未达到4\% & \textsf{\seqsplit{T$\in$[L,+∞)}}，其中 \textsf{\seqsplit{L}} 为最后一次检测孕周 & 7 \\ \hline
\end{longtable}
\endgroup

执行要求：三类样本全部保留；不使用区间中点代替首次达标时间，不把第一次达标检测孕周视为精确事件时间。

\subsection{规则 DC-04：缺失值分类处理}
\begingroup
\footnotesize
\setlength{\tabcolsep}{1.5pt}
\begin{longtable}{|>{\raggedright\arraybackslash}p{0.420\textwidth}|>{\raggedright\arraybackslash}p{0.420\textwidth}|}
\hline
缺失类型 & 当前处理 \\ \hline
\endfirsthead
\hline
缺失类型 & 当前处理 \\ \hline
\endhead
男胎末次月经缺失12条 & 问题1—3直接使用检测孕周，不使用末次月经；缺失值保留但不进入模型。 \\ \hline
男胎非整倍体字段空白956条 & 该字段不作为问题1—3的输入或结局，不进行填补。 \\ \hline
问题1—3关键建模字段缺失 & 不做均值、众数等隐式填补；缺失、非有限或无法转换时停止运行并检查数据。 \\ \hline
女胎两列全空无名列 & 直接删除，不进入特征集合。 \\ \hline
女胎BMI缺失1条 & 在每一个外层/内层训练折内用该训练折BMI中位数填补；同时增加“BMI缺失”指示变量。验证折和测试折只使用对应训练折中位数。 \\ \hline
女胎非整倍体字段空白 & 按附件输出语义解释为未列出T13、T18、T21异常，相应标签记为0；该标签不等同于真实临床诊断终点。 \\ \hline
\end{longtable}
\endgroup

执行要求：所有填补参数只能从当前训练折计算，禁止使用全数据或测试折统计量；除上述已定义情形外，不自动填补未知字段。

\section{三、第一问解答}
\subsection{1. 数据特征分析}
\begingroup
\footnotesize
\setlength{\tabcolsep}{1.5pt}
\begin{longtable}{|>{\raggedright\arraybackslash}p{0.280\textwidth}|>{\raggedright\arraybackslash}p{0.280\textwidth}|>{\raggedright\arraybackslash}p{0.280\textwidth}|}
\hline
项目 & 数据特征 & 对建模的影响 \\ \hline
\endfirsthead
\hline
项目 & 数据特征 & 对建模的影响 \\ \hline
\endhead
样本结构 & 男胎1082条检测记录对应267名孕妇，同一孕妇存在重复检测 & 记录之间并非相互独立，需要显式处理孕妇内相关性 \\ \hline
检测孕周 & 范围11—29周，均值约16.846周 & 可以分析Y浓度随孕周的纵向变化，并需要考虑非线性 \\ \hline
BMI & 范围20.703—46.875，均值约32.289 & 样本整体BMI偏高，需检验BMI主效应及其与孕周的交互 \\ \hline
Y染色体浓度 & 均值0.07719，中位数0.07507，为0—1之间的比例变量 & 不宜直接无限外推，需要进行logit变换 \\ \hline
其他变量 & 包括年龄、身高、IVF、读段数、比对比例和GC含量 & 可作为扩展变量检验是否提供额外统计解释量 \\ \hline
\end{longtable}
\endgroup

第一问的核心任务是解释变量关系与显著性，而不是寻找最佳检测时点或构建分类器。

\subsection{2. 数据预处理方案}
\begingroup
\footnotesize
\setlength{\tabcolsep}{1.5pt}
\begin{longtable}{|>{\raggedright\arraybackslash}p{0.280\textwidth}|>{\raggedright\arraybackslash}p{0.280\textwidth}|>{\raggedright\arraybackslash}p{0.280\textwidth}|}
\hline
处理步骤 & 具体做法 & 目的 \\ \hline
\endfirsthead
\hline
处理步骤 & 具体做法 & 目的 \\ \hline
\endhead
数据身份核验 & 检查附件SHA-256、1082条记录、267名孕妇及首末序号 & 防止输入版本发生变化 \\ \hline
孕周转换 & 将\textsf{\seqsplit{w+d}}转换为\textsf{\seqsplit{w+d/7}}十进制周 & 支持数值计算、排序和曲线建模 \\ \hline
重复检测 & 问题1保留全部1082条记录，以孕妇代码标识重复测量 & 保留Y浓度随孕周变化的纵向信息 \\ \hline
Y浓度变换 & 将Y浓度截断保护到\textsf{\seqsplit{[10\textasciicircum{}-6,1-10\textasciicircum{}-6]}}后进行logit变换 & 适应比例变量边界并避免0、1导致无穷值 \\ \hline
连续变量标准化 & 孕周、BMI、年龄、身高、比对比例和GC含量减均值后除以标准差 & 改善模型数值稳定性，并使系数可按一个标准差解释 \\ \hline
读段数处理 & 唯一比对读段数先取自然对数，再标准化 & 缓解量级过大和右偏问题 \\ \hline
IVF编码 & 自然受孕记为0，其他IVF方式记为1 & 将类别变量转换为模型可用的二元变量 \\ \hline
缺失检查 & Y浓度等关键字段缺失时停止运行，不做隐式填补 & 避免用未经确认的填补假设影响显著性 \\ \hline
\end{longtable}
\endgroup

\subsection{3. 模型选取及缘由}
主模型选用以孕妇为随机截距的线性混合效应模型，原因如下：

\begin{enumerate}
\item 同一孕妇的多次检测存在相关性，普通线性回归的独立性假设不成立。
\item 随机截距可以描述不同孕妇Y浓度基础水平的个体差异。
\item 固定效应系数、置信区间和显著性检验便于直接回答题目要求的关系方向。
\item 加入孕周二次项可以描述非线性变化；加入孕周×BMI交互项可以检验BMI是否改变孕周效应。
\item 第一问属于统计关系解释，不以搜索树或黑箱预测模型作为主模型。
\end{enumerate}

\subsection{4. 模型及其参数介绍}
基础模型为

\[
\operatorname{logit}(Y_{ij})=\beta_0+\beta_1W_{ij}+\beta_2W_{ij}^{2}
+\beta_3B_{ij}+\beta_4W_{ij}B_{ij}+u_i+\varepsilon_{ij},
\]
其中：

\begingroup
\footnotesize
\setlength{\tabcolsep}{1.5pt}
\begin{longtable}{|>{\raggedright\arraybackslash}p{0.420\textwidth}|>{\raggedright\arraybackslash}p{0.420\textwidth}|}
\hline
参数 & 含义 \\ \hline
\endfirsthead
\hline
参数 & 含义 \\ \hline
\endhead
\textsf{\seqsplit{$\beta$0}} & 所有标准化连续变量处于均值附近时的总体截距 \\ \hline
\textsf{\seqsplit{$\beta$1}} & 平均孕周附近，孕周增加一个标准差对应的Y浓度logit变化 \\ \hline
\textsf{\seqsplit{$\beta$2}} & 孕周关系的曲率，用于判断是否偏离直线关系 \\ \hline
\textsf{\seqsplit{$\beta$3}} & 相同孕周附近，BMI增加一个标准差对应的Y浓度logit变化 \\ \hline
\textsf{\seqsplit{$\beta$4}} & 孕周与BMI交互效应，检验BMI是否改变孕周增长斜率 \\ \hline
\textsf{\seqsplit{u\_i}} & 第\textsf{\seqsplit{i}}名孕妇的随机截距，\textsf{\seqsplit{u\_i\textasciitilde{}N(0,$\sigma$\_u²)}} \\ \hline
\textsf{\seqsplit{$\varepsilon$\_ij}} & 第\textsf{\seqsplit{i}}名孕妇第\textsf{\seqsplit{j}}次检测的记录级误差，\textsf{\seqsplit{$\varepsilon$\_ij\textasciitilde{}N(0,$\sigma$²)}} \\ \hline
\end{longtable}
\endgroup

扩展模型在基础模型上增加年龄、身高、IVF、对数唯一比对读段数、参考基因组比对比例和GC含量。基础模型和扩展模型均采用最大似然估计，优化方法为L-BFGS，最多迭代1000次。

\subsection{5. 模型训练}
第一问中的\textsf{\seqsplit{$\beta$}}、随机截距方差和残差方差不是人工设定或网格搜索得到的，而是由全部1082条纵向记录通过最大似然估计共同求得。

\begin{enumerate}
\item 先按照给定参数计算每条记录的条件均值，并结合孕妇随机截距构造整组纵向观测的联合似然。
\item 对随机效应积分后得到边际对数似然；代码使用L-BFGS迭代寻找使对数似然最大的参数组合。
\item 基础模型和扩展模型均使用\textsf{\seqsplit{REML=False}}的最大似然估计，确保两模型的似然值可以用于似然比检验。
\item 迭代最多1000次；若模型未收敛，则停止接受当前结果，不能直接引用系数。
\item 参数标准误由拟合点附近的曲率信息估计；Wald统计量、p值及95\%置信区间均由“系数÷标准误”及其近似分布得到。
\end{enumerate}

这里不存在“训练集上调参、测试集上评分”的常规机器学习流程，因为第一问主要进行统计关系估计。若未来增加预测性验证，应按孕妇分组另做交叉验证，不能改变当前显著性检验口径。

\subsection{6. 模型评分方式}
\begingroup
\footnotesize
\setlength{\tabcolsep}{1.5pt}
\begin{longtable}{|>{\raggedright\arraybackslash}p{0.280\textwidth}|>{\raggedright\arraybackslash}p{0.280\textwidth}|>{\raggedright\arraybackslash}p{0.280\textwidth}|}
\hline
评价方法 & 作用 & 判定方式 \\ \hline
\endfirsthead
\hline
评价方法 & 作用 & 判定方式 \\ \hline
\endhead
Wald检验 & 判断单个固定效应系数是否显著偏离0 & 结合系数、标准误、p值和95\%置信区间判断 \\ \hline
AIC & 比较模型拟合能力与复杂度的平衡 & AIC越低通常越优 \\ \hline
似然比检验 & 检验扩展变量作为一个整体是否增加解释量 & 比较基础模型与扩展模型对数似然，报告LR、自由度和p值 \\ \hline
收敛检查 & 确认参数优化是否成功 & 任一模型未收敛则不接受结果 \\ \hline
随机效应与残差 & 判断孕妇间差异及记录级误差规模 & 比较随机截距标准差和残差标准差 \\ \hline
\end{longtable}
\endgroup

第一问以参数估计和统计推断为目标，因此当前不使用训练集准确率作为评分，也不把预测交叉验证作为显著性结论的主要依据。

\subsection{7. 数值结果分析}
\begingroup
\footnotesize
\setlength{\tabcolsep}{1.5pt}
\begin{longtable}{|>{\raggedright\arraybackslash}p{0.280\textwidth}|>{\raggedright\arraybackslash}p{0.280\textwidth}|>{\raggedright\arraybackslash}p{0.280\textwidth}|}
\hline
变量/检验 & 结果 & 解释 \\ \hline
\endfirsthead
\hline
变量/检验 & 结果 & 解释 \\ \hline
\endhead
孕周一次项 & \textsf{\seqsplit{$\beta$=0.1668，p=3.32×10\textasciicircum{}-48}} & 孕周与Y浓度显著正相关 \\ \hline
孕周二次项 & \textsf{\seqsplit{$\beta$=0.0428，p=5.80×10\textasciicircum{}-6}} & 孕周关系存在显著曲率，不能完全视为恒定斜率 \\ \hline
BMI主效应 & \textsf{\seqsplit{$\beta$=-0.0719，p=0.00250}} & 相近孕周下，BMI越高，Y浓度整体越低 \\ \hline
孕周×BMI & \textsf{\seqsplit{$\beta$=0.00843，p=0.336}} & 交互不显著，不能声称不同BMI人群的增长速度显著不同 \\ \hline
基础模型AIC & \textsf{\seqsplit{836.330}} & 基础模型比较基准 \\ \hline
扩展模型AIC & \textsf{\seqsplit{817.693}} & 加入辅助变量后AIC降低18.637 \\ \hline
似然比检验 & \textsf{\seqsplit{LR=30.637，df=6，p=2.97×10\textasciicircum{}-5}} & 年龄、身高、IVF及测序变量整体增加统计解释量 \\ \hline
孕妇随机截距标准差 & \textsf{\seqsplit{0.434}} & 孕妇之间的基础Y浓度差异明显 \\ \hline
记录级残差标准差 & \textsf{\seqsplit{0.261}} & 同一孕妇内部仍存在检测波动和未解释误差 \\ \hline
\end{longtable}
\endgroup

\subsection{8. 结论解释与边界}
第一问可以得出：Y染色体浓度随孕周总体升高，并存在一定非线性；BMI与Y浓度显著负相关。孕周与BMI交互项不显著，因此当前证据只支持“高BMI孕妇的Y浓度整体偏低”，不支持“不同BMI人群具有显著不同的增长速度”。扩展变量整体改善模型解释量，但不能把单项统计系数直接解释为因果作用。

该结论来自当前267名孕妇的内部纵向数据。样本来源单一且BMI分布偏高，因此不能直接外推到所有地区或临床人群。第一问的结果为问题2、3按BMI进行时点分流提供方向依据，但不能单独决定最终BMI切点和推荐孕周。

\subsection{9. 补充}
\begin{enumerate}
\item 若混合模型不收敛、随机效应退化或残差出现非有限值，应停止使用该结果，并考虑患者聚类稳健GEE作为同一解释目标下的备选实现。
\item 回退模型与混合模型的参数含义不同，必须在报告中明确区分，不能混合引用系数。
\item 如需评价第一问模型的预测外推能力，可另做按孕妇分组的交叉验证；该结果只能作为补充，不能替代Wald检验和似然比检验。
\item 所有系数均作用于标准化变量和Y浓度logit尺度，不能直接解释为Y浓度增加或减少相同数值的百分点。
\end{enumerate}

\section{四、第二问解答}
\subsection{1. 数据特征分析}
本部分只描述数据现象如何阻碍“确定BMI分组及最佳NIPT时点”这一目标。

\begingroup
\footnotesize
\setlength{\tabcolsep}{1.5pt}
\begin{longtable}{|>{\raggedright\arraybackslash}p{0.420\textwidth}|>{\raggedright\arraybackslash}p{0.420\textwidth}|}
\hline
数据现象 & 对问题2的阻碍 \\ \hline
\endfirsthead
\hline
数据现象 & 对问题2的阻碍 \\ \hline
\endhead
1082条检测记录对应267名孕妇 & 检测记录不独立，直接按记录训练或划分会重复计权并产生患者泄漏 \\ \hline
同一孕妇仅在若干离散孕周接受检测 & 无法观察首次达到4\%的精确时点，只能形成左删失、区间删失或右删失 \\ \hline
左删失217人、区间删失43人、右删失7人 & 普通回归无法正确使用三类不完整时间信息 \\ \hline
BMI范围20.703—46.875且分布不均衡 & 高BMI尾部样本较少，切点可能受少量样本扰动，不能机械追求等人数分组 \\ \hline
达标概率通常随孕周增加 & 越晚检测越容易达到4\%，但更晚发现潜在问题；“准确”与“尽早”存在冲突 \\ \hline
题目把孕期风险分为$\leq$12周、13—27周、$\geq$28周三级 & 18周和22周在题面中同属中期；是否认为18周额外优于22周取决于风险解释 \\ \hline
同孕妇同孕周重复测量存在波动 & 4\%附近的状态可能因测量误差改变，需要检查推荐方案是否对扰动敏感 \\ \hline
\end{longtable}
\endgroup

\subsection{2. 问题分析}
第二问需要依次解决四类问题：

\begingroup
\footnotesize
\setlength{\tabcolsep}{1.5pt}
\begin{longtable}{|>{\raggedright\arraybackslash}p{0.420\textwidth}|>{\raggedright\arraybackslash}p{0.420\textwidth}|}
\hline
大类问题 & 需要回答的具体内容 \\ \hline
\endfirsthead
\hline
大类问题 & 需要回答的具体内容 \\ \hline
\endhead
首次达标时间预测 & 如何利用三类删失信息估计孕妇在不同孕周前达到4\%的概率 \\ \hline
BMI分组 & 分几组、切点设在哪里、如何避免局部贪心和极小样本组 \\ \hline
分组稳定孕周 & 每个BMI组在哪个最早整数周能够稳定满足达标可靠性要求 \\ \hline
题目风险解读 & “越早越好”与“通过率越高越好”如何同时进入判优；题面三级风险和连续等待风险如何区分 \\ \hline
\end{longtable}
\endgroup

这里“推荐周期越早越好”是因为检测推迟可能延后后续确认与处理；但时间早不能凌驾于检测可靠性之上。正确的判优逻辑是：先满足可靠性和稳定性，再在安全候选中选择更早、结构更简单的方案。

\subsection{3. 指标定义}
设第\textsf{\seqsplit{i}}名孕妇推荐孕周为\textsf{\seqsplit{t\_i}}，模型预测其在该周前达到4\%的概率为\textsf{\seqsplit{F\_i(t\_i)}}。

\begingroup
\footnotesize
\setlength{\tabcolsep}{1.5pt}
\begin{longtable}{|>{\raggedright\arraybackslash}p{0.280\textwidth}|>{\raggedright\arraybackslash}p{0.280\textwidth}|>{\raggedright\arraybackslash}p{0.280\textwidth}|}
\hline
指标 & 定义或门槛 & 含义 \\ \hline
\endfirsthead
\hline
指标 & 定义或门槛 & 含义 \\ \hline
\endhead
总体预测达标概率 & \textsf{\seqsplit{mean[F\_i(t\_i)]}} & 推荐时点前达到4\%的平均预测概率 \\ \hline
准确性风险 & \textsf{\seqsplit{R\_a=1-mean[F\_i(t\_i)]}} & 推荐时点仍未达到4\%的平均风险 \\ \hline
平均推荐孕周 & \textsf{\seqsplit{mean(t\_i)}} & 整体检测时间成本，越小表示越早 \\ \hline
题面分段时间风险 & 按$\leq$12周、13—27周、$\geq$28周分级 & 严格按照题目风险区间比较等待损失 \\ \hline
连续时间风险 & \textsf{\seqsplit{R\_t=mean[clip((t\_i-13)/14,0,1)]}} & 敏感性分析中假设13—27周内部等待损失连续增加 \\ \hline
不稳定风险 & \textsf{\seqsplit{R\_s=1-重复折安全率}} & 方案对重新划分数据的敏感程度 \\ \hline
综合风险 & \textsf{\seqsplit{R=w\_aR\_a+w\_tR\_t+w\_sR\_s}}，权重非负且和为1 & 比较不同准确性、等待和稳定性偏好下的方案 \\ \hline
组内平均达标安全线 & 至少0.89 & 固定政策部署审计的组内可靠性要求 \\ \hline
Bootstrap安全率 & 至少0.95 & 患者重采样后仍满足安全线的比例 \\ \hline
重复折安全率 & 至少0.90 & 重复划分中安全折所占比例 \\ \hline
重复4/4比例 & 至少0.70 & 100次重复四折中四折全部安全的比例 \\ \hline
\end{longtable}
\endgroup

所有门槛在测试审计前固定，不能看到测试结果后临时降低。

\subsection{4. 模型与算法组件及选取缘由}
\begingroup
\footnotesize
\setlength{\tabcolsep}{1.5pt}
\begin{longtable}{|>{\raggedright\arraybackslash}p{0.210\textwidth}|>{\raggedright\arraybackslash}p{0.210\textwidth}|>{\raggedright\arraybackslash}p{0.210\textwidth}|>{\raggedright\arraybackslash}p{0.210\textwidth}|}
\hline
模型/算法 & 是什么 & 选取缘由 & 解决的问题 \\ \hline
\endfirsthead
\hline
模型/算法 & 是什么 & 选取缘由 & 解决的问题 \\ \hline
\endhead
区间删失XGBoost-AFT & 将首次达到4\%视为生存时间，用提升树学习特征与达标时间分布的非线性关系 & 能直接接收左删失、区间删失和右删失上下界，并表示BMI等变量的非线性和交互 & 首次达标时间不精确、普通回归无法使用删失样本的问题 \\ \hline
患者级嵌套交叉验证 & 外层测试泛化、内层调参和早停，全部按孕妇划分 & 防止同一孕妇跨集合和测试折参与调参 & 患者泄漏、过拟合和不同方案比较不公平的问题 \\ \hline
全局Optimal Binning & 按BMI排序，只在相邻不同BMI之间建立连续分段，并比较完整候选分组 & 相比单棵决策树的局部贪心切分，更适合寻找整体最优、可解释的连续BMI区间 & BMI分几组、切点在哪里及局部最优问题 \\ \hline
组内最早安全周搜索 & 对每个候选组逐周计算预测达标概率并寻找满足安全要求的最早周 & 把连续概率预测转换为可执行的整数孕周建议 & 每组何时检测的问题 \\ \hline
Pareto/固定方案比较 & 同时比较达标概率、平均孕周、稳定性和复杂度 & 避免只追求最早或只追求最高通过率 & 多目标判优问题 \\ \hline
患者Bootstrap、重复折与误差扰动 & 重采样患者、重复划分并扰动4\%附近Y浓度 & 检查切点和时点是否由偶然样本或检测误差决定 & 稳定性与测量误差问题 \\ \hline
\end{longtable}
\endgroup

XGBoost-AFT只负责产生\textsf{\seqsplit{F\_i(t)=P(T\_i$\leq$t|x\_i)}}；Optimal Binning负责BMI分组；安全周搜索负责推荐孕周。三者不能混写成一个单独算法。

\subsection{5. 整体路线/算法结构}
\begin{quote}\small\sffamily
1082条男胎检测记录\\
        ↓ 同孕妇同孕周Y浓度取中位数、按孕周排序\\
267条患者级删失记录\\
        ↓ 外4折、内3折患者级嵌套验证\\
区间删失XGBoost-AFT\\
        ↓ 得到不同孕妇、不同孕周的OOF达标概率\\
BMI稳定排序＋合法相邻切点\\
        ↓ 全局Optimal Binning比较连续分组\\
候选BMI组\\
        ↓ 逐组搜索满足可靠性的最早整数周\\
候选分组—时点政策\\
        ↓ 固定比较15/18、16/18、18/22\\
题面风险、连续风险、Bootstrap、重复折和误差审计\\
        ↓\\
题面主推荐18/22＋连续风险敏感性方案16/18\\
\end{quote}

\subsection{6. 训练方案}
第二问包含三次相互衔接但目标不同的训练/优化：第一次训练达标时间分布，第二次训练BMI分组政策，第三次训练各组推荐时间。后两次本质上是基于冻结预测结果的决策参数优化，不再重新拟合Y浓度预测模型。

\subsubsection{6.1 第一次训练：XGBoost-AFT删失建模}
\textbf{训练目标}：不猜测精确达标时间，直接利用左删失、区间删失和右删失上下界，学习每名孕妇首次达到4\%的条件时间分布。

输入患者级数据样貌：

\begingroup
\footnotesize
\setlength{\tabcolsep}{1.5pt}
\begin{longtable}{|>{\raggedright\arraybackslash}p{0.280\textwidth}|>{\raggedright\arraybackslash}p{0.280\textwidth}|>{\raggedright\arraybackslash}p{0.280\textwidth}|}
\hline
字段 & 含义 & 示例 \\ \hline
\endfirsthead
\hline
字段 & 含义 & 示例 \\ \hline
\endhead
孕妇代码 & 患者唯一标识 & \textsf{\seqsplit{A001}} \\ \hline
BMI等基线特征 & 检测前可以获得的特征 & BMI=28.516 \\ \hline
\textsf{\seqsplit{lower}} & 首次达标时间下界 & 15.857周 \\ \hline
\textsf{\seqsplit{upper}} & 首次达标时间上界 & 20.143周或\textsf{\seqsplit{+∞}} \\ \hline
\textsf{\seqsplit{censoring}} & \textsf{\seqsplit{left}}、\textsf{\seqsplit{interval}}或\textsf{\seqsplit{right}} & \textsf{\seqsplit{interval}} \\ \hline
\textsf{\seqsplit{outer\_fold}} & 固定患者级外层折号 & 0—3 \\ \hline
\end{longtable}
\endgroup

训练过程：

\begin{enumerate}
\item 由1082条记录构造267条患者级删失记录，按删失类型分层。
\item 外层4折中，每次约200名孕妇训练、66—67名孕妇测试；外层训练集内部固定3折用于超参数选择和早停。
\item 对左删失、区间删失和右删失样本分别最大化\textsf{\seqsplit{F(U|x)}}、\textsf{\seqsplit{F(U|x)-F(L|x)}}和\textsf{\seqsplit{1-F(L|x)}}，等价地最小化其负对数似然。
\item AFT误差分布、尺度及树模型参数只使用当前外层训练数据确定；当前四个外层折均选择极值分布，尺度约为0.733、0.781、0.778和0.578。
\item 每个外层模型只预测该折未见过的孕妇，最后拼接成完整OOF概率表。
\end{enumerate}

XGBoost-AFT并没有把删失样本“补成一个确定孕周”。它把每个删失区间转换为一条随孕周单调上升的概率曲线。

第一次训练后的数据样貌：

\begingroup
\footnotesize
\setlength{\tabcolsep}{1.5pt}
\begin{longtable}{|>{\raggedright\arraybackslash}p{0.420\textwidth}|>{\raggedright\arraybackslash}p{0.420\textwidth}|}
\hline
数据规模/字段 & 结果 \\ \hline
\endfirsthead
\hline
数据规模/字段 & 结果 \\ \hline
\endhead
行数 & 267行，每名孕妇一行 \\ \hline
身份与标签字段 & 孕妇代码、BMI、\textsf{\seqsplit{lower}}、\textsf{\seqsplit{upper}}、删失类型、外层折号 \\ \hline
概率时间网格 & 10.000—25.000周，每隔1/7周计算一次 \\ \hline
概率字段数 & 106列，记为\textsf{\seqsplit{p\_10.000}}至\textsf{\seqsplit{p\_25.000}} \\ \hline
单元格含义 & \textsf{\seqsplit{p\_t=P(T$\leq$t|x)}}，表示在\textsf{\seqsplit{t}}周前达到4\%的OOF预测概率 \\ \hline
\end{longtable}
\endgroup

例如A001的删失区间为\textsf{\seqsplit{[15.857,20.143]}}，其OOF达标概率从10周的0.6936上升到16周的0.9305、18周的0.9620和22周的0.9902。该概率表是后两次训练共同使用的冻结输入。

\subsubsection{6.2 第二次训练：全局Optimal Binning确定BMI组数与切点}
\textbf{训练目标}：在BMI连续、有序且每组样本量可靠的约束下，确定分几组以及切点在哪里，使各组存在足够早且安全的推荐孕周。

最终采用全局有序分组搜索：

\begin{enumerate}
\item 将267名孕妇按BMI稳定排序，只允许在相邻不同BMI值之间切分。
\item 每个开发折搜索\textsf{\seqsplit{1..floor(n开发/30)}}个连续组；约200名开发患者时覆盖1—6组，全样本理论上最多8组。
\item 对每个候选区间读取AFT的OOF周概率，寻找单侧90\%置信下限达到0.95的最早整数孕周。
\item 在满足组规模和可靠性约束的完整分组中，通过动态规划最小化患者加权平均推荐孕周。
\item 在最优开发平均孕周0.25周以内优先选择组数更少的结构，再到完全留出的外层测试患者评价。
\end{enumerate}

四个外折中一组和两组各被选择2次，3组及以上均未胜出；平均测试孕周18.949周、平均测试达标概率0.9556，均值安全4/4折、保守LCB安全3/4折。因此数据支持简单分组结构，最终为部署可解释性采用两组政策。

第二次训练的最终分组：

\begingroup
\footnotesize
\setlength{\tabcolsep}{1.5pt}
\begin{longtable}{|>{\raggedright\arraybackslash}p{0.280\textwidth}|>{\raggedright\arraybackslash}p{0.280\textwidth}|>{\raggedright\arraybackslash}p{0.280\textwidth}|}
\hline
最终BMI范围 & 人数 & 占比 \\ \hline
\endfirsthead
\hline
最终BMI范围 & 人数 & 占比 \\ \hline
\endhead
\textsf{\seqsplit{BMI<34.357}} & 215 & 80.5\% \\ \hline
\textsf{\seqsplit{BMI$\geq$34.357}} & 52 & 19.5\% \\ \hline
\end{longtable}
\endgroup

固定切点网格审计表明，在每组至少30人并完成4个外折的条件下，29.75—34.50为均值与LCB均4/4安全的平台；算法切点34.357位于该平台内。

\subsubsection{6.3 第三次训练：最优推荐时间与安全裕量校准}
\textbf{训练目标}：在第二次训练固定的两组结构下，为低BMI组和高BMI组选择可执行的推荐孕周，处理“越早越好”和“达标概率越高越好”的冲突。

使用100次重复四折和患者Bootstrap比较最终可执行方案：

\begingroup
\scriptsize
\setlength{\tabcolsep}{1.5pt}
\begin{longtable}{|>{\raggedright\arraybackslash}p{0.168\textwidth}|>{\raggedright\arraybackslash}p{0.168\textwidth}|>{\raggedright\arraybackslash}p{0.168\textwidth}|>{\raggedright\arraybackslash}p{0.168\textwidth}|>{\raggedright\arraybackslash}p{0.168\textwidth}|}
\hline
最终比较方案 & Bootstrap安全率 & 重复折安全率 & 重复4/4比例 & 定位 \\ \hline
\endfirsthead
\hline
最终比较方案 & Bootstrap安全率 & 重复折安全率 & 重复4/4比例 & 定位 \\ \hline
\endhead
15/18 & 0.9812 & 0.83 & 0.41 & 重复划分不稳定，淘汰 \\ \hline
16/18 & 0.9976 & 0.94 & 0.77 & 连续等待风险下的效率折中方案 \\ \hline
18/22 & 1.0000 & 1.00 & 1.00 & 题面三级风险下的主推荐 \\ \hline
\end{longtable}
\endgroup

第三次训练最终没有只输出一个“数学最早周”，而是输出两层政策：题面把18周和22周都视为中期，因此主推荐选择准确性和稳定性更高的18/22；如果额外假定13—27周内部等待损失连续增加，则保留16/18作为风险权重敏感性方案。

\subsection{7. 验证方案}
\begingroup
\footnotesize
\setlength{\tabcolsep}{1.5pt}
\begin{longtable}{|>{\raggedright\arraybackslash}p{0.280\textwidth}|>{\raggedright\arraybackslash}p{0.280\textwidth}|>{\raggedright\arraybackslash}p{0.280\textwidth}|}
\hline
验证层 & 做法 & 作用 \\ \hline
\endfirsthead
\hline
验证层 & 做法 & 作用 \\ \hline
\endhead
外层4折 & 按删失类型分层，每名孕妇仅出现于一个测试折 & 生成未参与该折训练的患者级OOF概率 \\ \hline
内层3折 & 仅在外层训练集内调参和早停 & 隔离模型选择与外层评价 \\ \hline
固定政策回放 & 在同一批冻结预测上比较15/18、16/18和18/22 & 保证不同政策使用相同预测基础 \\ \hline
患者Bootstrap & 以孕妇为单位重采样，不独立抽取检测记录 & 估计组内可靠性和方案排序的抽样波动 \\ \hline
100次重复四折 & 改变患者折分后重新检查各折安全性 & 评价方案对数据划分的稳定性 \\ \hline
测量误差敏感性 & 根据同孕妇同孕周重复检测估计误差并扰动Y浓度 & 检查4\%附近状态改变是否推翻方案排序 \\ \hline
\end{longtable}
\endgroup

上述属于内部验证，不是来自新医院或新时间段的外部测试。最终切点和固定政策经过当前数据开发，不能把内部回放结果表述为完全独立的无偏临床成绩。

\subsection{8. 结论对比}
\begingroup
\tiny
\setlength{\tabcolsep}{1.5pt}
\begin{longtable}{|>{\raggedright\arraybackslash}p{0.120\textwidth}|>{\raggedright\arraybackslash}p{0.120\textwidth}|>{\raggedright\arraybackslash}p{0.120\textwidth}|>{\raggedright\arraybackslash}p{0.120\textwidth}|>{\raggedright\arraybackslash}p{0.120\textwidth}|>{\raggedright\arraybackslash}p{0.120\textwidth}|>{\raggedright\arraybackslash}p{0.120\textwidth}|}
\hline
方案 & 平均推荐孕周 & 总体达标概率 & Bootstrap安全率 & 重复折安全率 & 重复4/4比例 & 判定 \\ \hline
\endfirsthead
\hline
方案 & 平均推荐孕周 & 总体达标概率 & Bootstrap安全率 & 重复折安全率 & 重复4/4比例 & 判定 \\ \hline
\endhead
15/18 & 15.584 & 0.9102 & 0.9812 & 0.83 & 0.41 & 过早且重复折不稳定，淘汰 \\ \hline
16/18 & 16.390 & 0.9244 & 0.9976 & 0.94 & 0.77 & 满足内部稳定门槛，作为连续等待风险敏感性方案 \\ \hline
18/22 & 18.779 & 0.9555 & 1.0000 & 1.00 & 1.00 & 题面三级风险下准确性与稳定性最好，作为主推荐 \\ \hline
\end{longtable}
\endgroup

15/18虽然最早，但重复折安全率和4/4比例未达门槛；16/18在效率与可靠性之间折中；18/22最晚，但三套方案均处于题目规定的中期风险区间，因此题面口径下优先选择其更高的达标概率和稳定性。

\subsection{9. 最终结果}
题面三级风险主推荐：

\begingroup
\footnotesize
\setlength{\tabcolsep}{1.5pt}
\begin{longtable}{|>{\raggedright\arraybackslash}p{0.210\textwidth}|>{\raggedright\arraybackslash}p{0.210\textwidth}|>{\raggedright\arraybackslash}p{0.210\textwidth}|>{\raggedright\arraybackslash}p{0.210\textwidth}|}
\hline
BMI范围 & 人数 & 推荐NIPT时点 & 适用口径 \\ \hline
\endfirsthead
\hline
BMI范围 & 人数 & 推荐NIPT时点 & 适用口径 \\ \hline
\endhead
\textsf{\seqsplit{BMI<34.357}} & 215 & 第18周 & 题面三级风险；准确性与稳定性优先 \\ \hline
\textsf{\seqsplit{BMI$\geq$34.357}} & 52 & 第22周 & 题面三级风险；准确性与稳定性优先 \\ \hline
\end{longtable}
\endgroup

连续等待风险敏感性方案：

\begingroup
\footnotesize
\setlength{\tabcolsep}{1.5pt}
\begin{longtable}{|>{\raggedright\arraybackslash}p{0.210\textwidth}|>{\raggedright\arraybackslash}p{0.210\textwidth}|>{\raggedright\arraybackslash}p{0.210\textwidth}|>{\raggedright\arraybackslash}p{0.210\textwidth}|}
\hline
BMI范围 & 人数 & 推荐NIPT时点 & 适用口径 \\ \hline
\endfirsthead
\hline
BMI范围 & 人数 & 推荐NIPT时点 & 适用口径 \\ \hline
\endhead
\textsf{\seqsplit{BMI<34.357}} & 215 & 第16周 & 假定13—27周内等待损失连续增加 \\ \hline
\textsf{\seqsplit{BMI$\geq$34.357}} & 52 & 第18周 & 假定13—27周内等待损失连续增加 \\ \hline
\end{longtable}
\endgroup

主推荐的总体预测达标概率为0.9555。这里的达标概率是Y染色体浓度达到4\%的概率，不是胎儿异常诊断准确率。

\subsection{10. 补充}
\begin{enumerate}
\item \textsf{\seqsplit{34.357}}是当前数据和目标函数下的算法切点，不是通用医学BMI诊断界值；报告中可写为34.36。
\item 每组至少30人的约束用于防止统计上极不稳定的小组，不要求不同BMI组人数相等；最终两组为215人和52人。
\item 当\textsf{\seqsplit{w\_s=0}}时，16/18仅在准确性权重约\textsf{\seqsplit{0.8022—0.8461}}范围内最优，因此它依赖风险偏好，不能替代题面主答案。
\item 结果来自单一数据源和内部验证，高BMI尾部样本有限，尚无独立医院数据验证。
\item 推荐用于数学建模和研究性检测分流，不构成临床诊断或处置意见。
\end{enumerate}

\section{五、第三问解答}
\subsection{1. 数据特征分析}
沿用前述267名孕妇的患者级删失数据、重复检测聚合、孕周转换、患者级划分和删失区间构造规则，不再重复展开。

新增使用年龄、身高、体重、IVF、孕次和产次等检测前因素。当前主要障碍为：变量增加后过拟合风险上升；BMI、身高和体重之间存在相关性；高BMI尾部样本较少；平均达标概率可能掩盖组内高风险个体；多因素概率最终仍需压缩成可执行的BMI分组；测量误差可能改变4\%附近样本的达标状态。

\subsection{2. 问题分析}
\begingroup
\footnotesize
\setlength{\tabcolsep}{1.5pt}
\begin{longtable}{|>{\raggedright\arraybackslash}p{0.420\textwidth}|>{\raggedright\arraybackslash}p{0.420\textwidth}|}
\hline
大类问题 & 具体任务 \\ \hline
\endfirsthead
\hline
大类问题 & 具体任务 \\ \hline
\endhead
多因素达标建模 & 综合年龄、身高、体重、BMI、IVF、孕次和产次估计不同孕周前达到4\%的概率 \\ \hline
分组结构优化 & 根据多因素概率确定合理的BMI组数及连续切点 \\ \hline
组内安全周期 & 为每个BMI组寻找满足可靠性要求的最早整数孕周 \\ \hline
组内差异控制 & 同时考虑平均达标概率、组内个体差异和样本量 \\ \hline
结果稳定性 & 检查组数、切点和时点是否随折分或患者重采样发生明显变化 \\ \hline
测量误差处理 & 检查Y浓度扰动是否推翻推荐政策 \\ \hline
部署简化 & 将不稳定的小数切点转化为便于解释和执行的方案 \\ \hline
\end{longtable}
\endgroup

目标不是机械增加组数，而是在综合相关因素后，寻找可靠性、检测时间、稳定性和可解释性之间的平衡。

\subsection{3. 指标定义}
\subsubsection{3.1 维持定义}
总体预测达标概率、准确性风险、平均推荐孕周、题面分段时间风险、连续时间风险、不稳定风险、综合风险、Bootstrap安全率、重复折安全率和测量误差状态改变率沿用前述定义。

\subsubsection{3.2 新增定义}
设某一BMI组在孕周\textsf{\seqsplit{t}}的平均预测达标概率为\textsf{\seqsplit{p̄\_g(t)}}、组内标准差为\textsf{\seqsplit{s\_g(t)}}、人数为\textsf{\seqsplit{n\_g}}，定义单侧保守下界：

\[
LCB_g(t)=\bar p_g(t)-1.645\frac{s_g(t)}{\sqrt{n_g}}.
\]
新增指标如下：

\begingroup
\footnotesize
\setlength{\tabcolsep}{1.5pt}
\begin{longtable}{|>{\raggedright\arraybackslash}p{0.280\textwidth}|>{\raggedright\arraybackslash}p{0.280\textwidth}|>{\raggedright\arraybackslash}p{0.280\textwidth}|}
\hline
指标 & 定义 & 作用 \\ \hline
\endfirsthead
\hline
指标 & 定义 & 作用 \\ \hline
\endhead
组内保守下界 & \textsf{\seqsplit{LCB\_g(t)}} & 同时反映组内均值、个体差异和样本量 \\ \hline
均值安全折数 & 外层测试折中最低平均达标概率达到安全线的折数 & 评价平均意义下的跨折可靠性 \\ \hline
保守下限安全折数 & 外层测试折中最低LCB达到安全线的折数 & 评价更严格的跨折可靠性 \\ \hline
切点Bootstrap区间 & 切点Bootstrap分布的2.5\%和97.5\%分位数 & 评价精确切点的抽样稳定性 \\ \hline
组数选择稳定率 & 某组数在Bootstrap中被选择的次数除以总次数 & 评价分组结构是否反复出现 \\ \hline
\end{longtable}
\endgroup

\subsubsection{3.3 变更定义}
安全周由主要依据组内平均达标概率，调整为优先依据保守下界：

\[
t_g^*=\min\{t:LCB_g(t)\ge q\}.
\]
开发安全裕量取\textsf{\seqsplit{q=0.95}}。判优顺序调整为：先保证候选可行和保守可靠，再比较均值可靠性与平均推荐孕周；性能接近时优先采用组数更少、切点更稳定且更容易解释的结构。

\subsection{4. 模型与算法组件及选取缘由}
\begingroup
\footnotesize
\setlength{\tabcolsep}{1.5pt}
\begin{longtable}{|>{\raggedright\arraybackslash}p{0.210\textwidth}|>{\raggedright\arraybackslash}p{0.210\textwidth}|>{\raggedright\arraybackslash}p{0.210\textwidth}|>{\raggedright\arraybackslash}p{0.210\textwidth}|}
\hline
模型或算法 & 是什么 & 选取缘由 & 解决的问题 \\ \hline
\endfirsthead
\hline
模型或算法 & 是什么 & 选取缘由 & 解决的问题 \\ \hline
\endhead
多因素XGBoost-AFT & 在删失时间模型中加入BMI、年龄、身高、体重、IVF、孕次和产次 & 能处理删失并表示多因素非线性关系 & 生成综合相关因素后的达标概率 \\ \hline
有序动态规划 & 按BMI排序，在连续BMI区间中组合可行分段 & 固定约束下可比较完整切点组合，避免局部贪心 & 确定连续BMI分组及切点 \\ \hline
保守下界安全周搜索 & 根据LCB确定组内最早安全孕周 & 防止平均值掩盖组内高风险个体 & 确定稳定检测周期 \\ \hline
嵌套政策选择 & 每个外层训练折内部重新选择分组和时点 & 防止外层测试结果参与选参 & 评价整套优化流程的泛化能力 \\ \hline
患者Bootstrap & 以孕妇为单位重采样并重新执行政策审计 & 检查切点和政策是否由少量患者决定 & 评价抽样稳定性 \\ \hline
测量误差模拟 & 按重复检测估计的误差扰动Y浓度 & 检查4\%附近状态变化 & 评价最终政策对测量误差的敏感程度 \\ \hline
\end{longtable}
\endgroup

\subsection{5. 整体路线/算法结构}
\begin{quote}\small\sffamily
患者级删失数据\\
        ↓ 加入年龄、身高、体重、BMI、IVF、孕次和产次\\
多因素XGBoost-AFT\\
        ↓ 生成不同孕周前达到4\%的OOF概率\\
BMI稳定排序\\
        ↓ 有序动态规划搜索连续BMI分组\\
保守下界安全周搜索\\
        ↓ 外层训练折内部选择分组、切点和时点\\
嵌套外层审计\\
        ↓ Bootstrap切点与政策稳定性分析\\
测量误差模拟\\
        ↓ 简化切点并形成可执行政策\\
\end{quote}

\subsection{6. 训练方案}
\subsubsection{6.1 多因素达标概率训练}
输入特征为：

\[
x_i=(BMI,Age,Height,Weight,IVF,Gravidity,Parity).
\]
输出为：

\[
F_i(t)=P(T_i\le t\mid x_i).
\]
删失似然、患者级外4折、内3折和OOF概率生成方式沿用前述方案。输出仍为267行患者级概率曲线，后续优化读取冻结概率，不重新拟合AFT模型。

\subsubsection{6.2 BMI分组结构训练}
将孕妇按BMI稳定排序，只允许在相邻不同BMI值之间设置切点。对每个候选BMI区间，计算不同孕周下的组内均值和LCB，寻找满足\textsf{\seqsplit{LCB$\geq$0.95}}的最早孕周；不存在安全孕周的区间记为不可行。

有序动态规划组合多个可行区间，并在给定实验约束下最小化：

\[
\sum_g n_gt_g,
\]
即最小化患者加权平均推荐孕周。每个开发折搜索\textsf{\seqsplit{1..floor(n开发/30)}}个组；约200名开发患者时实际覆盖1—6组。

\subsubsection{6.3 推荐时点与部署政策训练}
在每个外层训练折内完成分组、切点和时点选择；以开发平均孕周为效率指标，在相差不超过0.25周的近优方案中优先选择更简单的结构。冻结后应用到该折测试孕妇，外层测试结果不参与当前折的选参。

全样本有序分组得到切点34.357；固定切点网格的4/4安全平台为29.75—34.50，因此采用34.357作为最终数据切点。

\subsection{7. 验证方案}
患者级外层4折、内层3折、患者Bootstrap和测量误差模拟沿用前述验证原则。

无外层选参审计中，每个外层折只依据该折训练数据选择分组和时点，得到：平均测试孕周18.949、平均测试达标概率0.9556、均值安全4/4折、保守下限安全3/4折。

\subsection{8. 结论对比}
\begingroup
\tiny
\setlength{\tabcolsep}{1.5pt}
\begin{longtable}{|>{\raggedright\arraybackslash}p{0.084\textwidth}|>{\raggedright\arraybackslash}p{0.084\textwidth}|>{\raggedright\arraybackslash}p{0.084\textwidth}|>{\raggedright\arraybackslash}p{0.084\textwidth}|>{\raggedright\arraybackslash}p{0.084\textwidth}|>{\raggedright\arraybackslash}p{0.084\textwidth}|>{\raggedright\arraybackslash}p{0.084\textwidth}|>{\raggedright\arraybackslash}p{0.084\textwidth}|>{\raggedright\arraybackslash}p{0.084\textwidth}|>{\raggedright\arraybackslash}p{0.084\textwidth}|}
\hline
外层测试折 & 开发集选择组数 & 开发集切点 & 开发集推荐孕周 & 测试平均孕周 & 测试总体达标概率 & 测试审计最低均值 & 测试审计最低LCB & 均值安全 & LCB安全 \\ \hline
\endfirsthead
\hline
外层测试折 & 开发集选择组数 & 开发集切点 & 开发集推荐孕周 & 测试平均孕周 & 测试总体达标概率 & 测试审计最低均值 & 测试审计最低LCB & 均值安全 & LCB安全 \\ \hline
\endhead
1 & 1 & 无 & 19 & 19.000 & 0.9598 & 0.9128 & 0.8938 & 是 & 否 \\ \hline
2 & 2 & 35.171 & 19/23 & 19.478 & 0.9573 & 0.9523 & 0.9452 & 是 & 是 \\ \hline
3 & 1 & 无 & 19 & 19.000 & 0.9583 & 0.9434 & 0.9274 & 是 & 是 \\ \hline
4 & 2 & 32.454 & 17/20 & 18.318 & 0.9469 & 0.9386 & 0.9283 & 是 & 是 \\ \hline
汇总 & 1组和2组各2折 & — & — & 18.949 & 0.9556 & — & — & 4/4折 & 3/4折 \\ \hline
\end{longtable}
\endgroup

全部可行组数搜索中，一组和两组各被选择2/4折，3组及以上均未胜出。各折测试总体达标概率为0.9469—0.9598，四折均值安全，三个外折同时满足保守LCB安全。综合可靠性、平均推荐孕周和执行可解释性，最终采用两组政策；该结论表示简单结构具有稳定优势，不表示两组在每次数据划分中都是唯一最优。

\subsection{9. 最终结果}
\begingroup
\footnotesize
\setlength{\tabcolsep}{1.5pt}
\begin{longtable}{|>{\raggedright\arraybackslash}p{0.280\textwidth}|>{\raggedright\arraybackslash}p{0.280\textwidth}|>{\raggedright\arraybackslash}p{0.280\textwidth}|}
\hline
BMI范围 & 人数 & 推荐NIPT时点 \\ \hline
\endfirsthead
\hline
BMI范围 & 人数 & 推荐NIPT时点 \\ \hline
\endhead
\textsf{\seqsplit{BMI<34.357}} & 215 & 第18周 \\ \hline
\textsf{\seqsplit{BMI$\geq$34.357}} & 52 & 第22周 \\ \hline
\end{longtable}
\endgroup

\begingroup
\footnotesize
\setlength{\tabcolsep}{1.5pt}
\begin{longtable}{|>{\raggedright\arraybackslash}p{0.420\textwidth}|>{\raggedright\arraybackslash}p{0.420\textwidth}|}
\hline
指标 & 数值 \\ \hline
\endfirsthead
\hline
指标 & 数值 \\ \hline
\endhead
平均推荐孕周 & 18.949 \\ \hline
总体预测达标概率 & 0.9556 \\ \hline
均值安全折数 & 4/4 \\ \hline
保守LCB安全折数 & 3/4 \\ \hline
\end{longtable}
\endgroup

\subsection{10. 补充}
\begin{enumerate}
\item 算法切点34.357位于29.75—34.50的4/4安全网格平台内；它是当前数据与目标函数共同确定的切点，不是通用医学BMI诊断界值。
\item 测量误差标准差约为0.003716，最终政策下平均状态改变率约为2.72\%，没有推翻18/22政策。
\item 保守LCB为3/4折安全，因此论文中应将结果表述为内部稳健推荐，而非绝对安全结论。
\item 当前结果属于内部验证，尚无独立医院或新时间段数据验证。
\end{enumerate}

\section{六、第四问解答}
\subsection{1. 数据特征分析}
女胎数据包含605条检测记录、147名孕妇和31个原始字段，平均每名孕妇约4.12条记录，最多9条。标签来自附件AB列“染色体的非整倍体”，将组合字符串拆为三个非互斥二元标签。

\begingroup
\footnotesize
\setlength{\tabcolsep}{1.5pt}
\begin{longtable}{|>{\raggedright\arraybackslash}p{0.210\textwidth}|>{\raggedright\arraybackslash}p{0.210\textwidth}|>{\raggedright\arraybackslash}p{0.210\textwidth}|>{\raggedright\arraybackslash}p{0.210\textwidth}|}
\hline
标签 & 阳性记录数 & 阳性孕妇数 & 数据障碍 \\ \hline
\endfirsthead
\hline
标签 & 阳性记录数 & 阳性孕妇数 & 数据障碍 \\ \hline
\endhead
T13 & 23 & 18 & 阳性较少，指标波动较大 \\ \hline
T18 & 46 & 30 & 三类中样本相对最多，但仍明显不平衡 \\ \hline
T21 & 13 & 12 & 极少样本，最容易出现不稳定和过拟合 \\ \hline
任一异常 & 67 & 44 & 适合总体风险排序，但不能替代染色体特异判定 \\ \hline
未列出异常 & 538 & — & 占大多数，普通准确率会产生误导 \\ \hline
\end{longtable}
\endgroup

附件存在T13T18、T13T21和T18T21等组合标签，因此不能建成互斥多分类。AE列“胎儿是否健康”605条均为“是”，没有可用于监督学习的真实临床结局；当前只能复现AB列筛查标签，不能声称诊断真实胎儿疾病。

同一孕妇存在重复检测，若随机拆分记录会导致同一孕妇同时进入训练集和测试集。BMI缺失1条，两列无名列全部为空；标签严重不平衡，筛查任务需要优先控制漏检，同时承担较多假阳性和复测负担。

\subsection{2. 问题分析}
\begingroup
\footnotesize
\setlength{\tabcolsep}{1.5pt}
\begin{longtable}{|>{\raggedright\arraybackslash}p{0.420\textwidth}|>{\raggedright\arraybackslash}p{0.420\textwidth}|}
\hline
大类问题 & 具体任务 \\ \hline
\endfirsthead
\hline
大类问题 & 具体任务 \\ \hline
\endhead
多标签异常判定 & 同时输出T13、T18和T21概率及非互斥标签 \\ \hline
类别不平衡 & 防止模型只预测大量阴性记录而获得虚高准确率 \\ \hline
重复检测 & 保证同一孕妇全部记录始终位于同一数据折，并限制多次检测孕妇的总权重 \\ \hline
概率与阈值 & 将连续风险概率转化为染色体特异阳性/阴性判定 \\ \hline
概率校准 & 使模型分数尽可能具有概率解释 \\ \hline
不确定记录 & 对阈值附近结果设置复测区，而不是强制二分类 \\ \hline
测序质量 & 识别读段、比对比例、GC等质量极端记录 \\ \hline
方法对照 & 与固定Z值规则和训练型Z值规则公平比较 \\ \hline
结论边界 & 明确输出属于附件筛查标签复现和研究性分流 \\ \hline
\end{longtable}
\endgroup

\subsection{3. 指标定义}
\subsubsection{3.1 维持定义}
患者级外层测试、内层OOF选参、患者Bootstrap、输入哈希固定和测试折不参与参数选择等验证定义沿用前述原则。模型稳定性继续通过不同外层折和患者重采样评价，不按检测记录独立抽样。

\subsubsection{3.2 新增定义}
设\textsf{\seqsplit{TP、TN、FP、FN}}分别为真阳性、真阴性、假阳性和假阴性。

\begingroup
\footnotesize
\setlength{\tabcolsep}{1.5pt}
\begin{longtable}{|>{\raggedright\arraybackslash}p{0.280\textwidth}|>{\raggedright\arraybackslash}p{0.280\textwidth}|>{\raggedright\arraybackslash}p{0.280\textwidth}|}
\hline
指标 & 定义 & 作用 \\ \hline
\endfirsthead
\hline
指标 & 定义 & 作用 \\ \hline
\endhead
灵敏度 & \textsf{\seqsplit{TP/(TP+FN)}} & 衡量阳性标签被检出的比例，筛查任务优先关注 \\ \hline
特异度 & \textsf{\seqsplit{TN/(TN+FP)}} & 衡量阴性记录被正确排除的比例 \\ \hline
阳性预测值PPV & \textsf{\seqsplit{TP/(TP+FP)}} & 衡量模型阳性中附件标签阳性的比例 \\ \hline
阴性预测值NPV & \textsf{\seqsplit{TN/(TN+FN)}} & 衡量模型阴性中附件标签阴性的比例 \\ \hline
F1 & 精确率和召回率的调和平均 & 平衡漏检和误报 \\ \hline
F2 & 对召回率赋予更高权重的F分数 & 符合筛查阶段漏检代价较高的取向 \\ \hline
ROC-AUC & 灵敏度—假阳性率曲线下面积 & 衡量概率排序能力，但不单独判优 \\ \hline
PR-AUC & 精确率—召回率曲线下面积 & 类别不平衡下的主排序指标 \\ \hline
Brier分数 & \textsf{\seqsplit{mean[(p-y)\textasciicircum{}2]}} & 衡量概率误差，越小越好 \\ \hline
复测率 & 被标记为阈值附近或质量异常的记录比例 & 衡量不确定性处理带来的执行成本 \\ \hline
患者Bootstrap区间 & 以孕妇为单位有放回抽样得到的指标分位数区间 & 衡量小样本和稀有阳性造成的不确定性 \\ \hline
\end{longtable}
\endgroup

\subsubsection{3.3 变更定义}
任务输出由“首次达标时间和推荐孕周”变更为“染色体异常概率、多标签判定和复测状态”。主判优指标变更为宏平均PR-AUC及稀有标签灵敏度；阈值选择优先保证灵敏度不低于85\%，再最大化特异度，并以F2处理并列。普通准确率不作为主指标。

任一异常概率采用三个染色体概率并集：

\[
p_i(any)=1-\prod_{c\in\{13,18,21\}}(1-p_{ic}).
\]
该概率用于总体风险排序；染色体特异性标签仍由三个独立阈值决定。

\subsection{4. 模型与算法组件及选取缘由}
\begingroup
\footnotesize
\setlength{\tabcolsep}{1.5pt}
\begin{longtable}{|>{\raggedright\arraybackslash}p{0.210\textwidth}|>{\raggedright\arraybackslash}p{0.210\textwidth}|>{\raggedright\arraybackslash}p{0.210\textwidth}|>{\raggedright\arraybackslash}p{0.210\textwidth}|}
\hline
模型或算法 & 是什么 & 选取缘由 & 解决的问题 \\ \hline
\endfirsthead
\hline
模型或算法 & 是什么 & 选取缘由 & 解决的问题 \\ \hline
\endhead
三个代价敏感XGBoost分类器 & 分别预测T13、T18和T21概率 & 能表示Z值、GC、读段质量和孕妇特征的非线性交互 & 多标签概率筛查 \\ \hline
患者等权 & 同一孕妇每条记录权重与其记录数倒数成正比 & 防止多次检测孕妇支配损失 & 重复检测计权偏差 \\ \hline
正类权重 & 每个训练折中以阴性数/阳性数提高正类损失权重 & 缓解阳性稀少导致的全阴性倾向 & 类别不平衡 \\ \hline
Platt校准 & 用内层OOF分数拟合单变量sigmoid概率映射 & 改善模型分数的概率解释 & 概率校准 \\ \hline
灵敏度优先阈值搜索 & 在内层OOF概率的全部候选阈值中按固定规则选择 & 阈值不读取外层测试标签 & 概率到判定的转换 \\ \hline
阈值不确定带 & 根据预测概率到阈值的距离确定复测区 & 避免对边界记录强制二分 & 不确定记录处理 \\ \hline
训练折质量分位数 & 根据读段、比对、重复、过滤和GC分布标记质量极端 & 不使用测试折统计量设定质量线 & 低质量检测识别 \\ \hline
Z值规则对照 & 固定\textsf{\seqsplit{|Z|$\geq$3}}及训练折内搜索的Z阈值 & 提供简单透明的传统基准 & 判断复杂模型是否提供额外筛查能力 \\ \hline
\end{longtable}
\endgroup

\subsection{5. 整体路线/算法结构}
\begin{quote}\small\sffamily
605条女胎检测记录\\
        ↓ 标签拆分、孕周转换、缺失和泄漏字段处理\\
T13/T18/T21三个非互斥标签\\
        ↓ 患者级外4折、内3折\\
三个代价敏感XGBoost分类器\\
        ↓ 内层OOF参数选择与Platt概率校准\\
染色体特异概率\\
        ↓ 灵敏度优先阈值搜索\\
T13/T18/T21多标签判定\\
        ↓ 概率并集形成任一异常风险\\
阈值不确定带＋测序质量规则\\
        ↓\\
筛查阳性／筛查阴性／阈值附近复测／质量异常重抽\\
        ↓\\
Z值规则对照＋患者Bootstrap稳定性评价\\
\end{quote}

\subsection{6. 训练方案}
\subsubsection{6.1 多标签概率模型训练}
对每个染色体\textsf{\seqsplit{c$\in$\{13,18,21\}}}分别训练分类器：

\[
p_{ic}=P(y_{ic}=1\mid x_i).
\]
输入包括目标染色体及X染色体Z值、X染色体浓度、读段数量与比例、GC含量、年龄、身高、体重、BMI、孕周、IVF、孕次、产次和抽血次数；同时构造绝对Z值、最大绝对Z值、Z值间差异、唯一比对读段占比、GC偏离量及BMI缺失指示。

每个训练折内计算患者等权\textsf{\seqsplit{v\_i}}和正类权重\textsf{\seqsplit{$\alpha$\_c}}，优化带权二元对数损失：

\[
L_c=-\sum_i v_i[\alpha_cy_{ic}\log p_{ic}+(1-y_{ic})\log(1-p_{ic})]+\Omega(f_c).
\]
内层候选配置比较树深、学习率、采样率、正则化和迭代轮数，主评分为PR-AUC，辅助参考F2和概率误差。

\subsubsection{6.2 概率校准与阈值训练}
用完整内层OOF原始概率拟合Platt校准器，再作用于外层测试原始概率。内层OOF阳性少于8条、阴性少于20条或校准器数值失败时，回退到原始XGBoost概率。

对每个染色体枚举相邻唯一概率的中点及两个端点作为候选阈值。阳性样本充足时，先保留灵敏度不低于85\%的阈值，再选择特异度最高者；并列时依次选择F2更高和阈值更高者。样本不足或无可行阈值时，改为优先最大化F2。

\subsubsection{6.3 复测与最终分流规则训练}
定义阈值距离\textsf{\seqsplit{d\_ic=|p\_ic-t\_c|}}，不确定带宽取训练OOF距离的15\%经验分位数且不超过0.15。若任一染色体位于不确定带，则建议复测。

测序质量规则只使用训练折分位数：六项质量指标中至少两项越界，或总/唯一读段数低于训练折0.5\%分位点时，标记质量异常并建议重抽复检。最终分流优先级为质量异常、阈值附近、筛查阳性、筛查阴性。

\subsection{7. 验证方案}
外层4折按孕妇分组，每折测试36—37名孕妇、151—152条记录；患者重叠数为0。内层3折完成特征处理、参数选择、校准、阈值和复测带宽确定。分折使用患者级多标签平衡搜索，使每个外层测试折均包含T13、T18和T21阳性孕妇。

每个染色体分别与固定\textsf{\seqsplit{|Z|$\geq$3}}规则和训练型Z阈值比较。患者级Bootstrap重复2000次，报告灵敏度、特异度、PR-AUC、ROC-AUC、F2和Brier分数的95\%区间。若稀有标签长期零召回、性能低于Z值基准、Bootstrap下限接近随机或复测负担失控，则不得宣称主模型有效。

\subsection{8. 结论对比}
\begingroup
\tiny
\setlength{\tabcolsep}{1.5pt}
\begin{longtable}{|>{\raggedright\arraybackslash}p{0.120\textwidth}|>{\raggedright\arraybackslash}p{0.120\textwidth}|>{\raggedright\arraybackslash}p{0.120\textwidth}|>{\raggedright\arraybackslash}p{0.120\textwidth}|>{\raggedright\arraybackslash}p{0.120\textwidth}|>{\raggedright\arraybackslash}p{0.120\textwidth}|>{\raggedright\arraybackslash}p{0.120\textwidth}|}
\hline
标签 & 方法 & 灵敏度 & 特异度 & PR-AUC & ROC-AUC & F2 \\ \hline
\endfirsthead
\hline
标签 & 方法 & 灵敏度 & 特异度 & PR-AUC & ROC-AUC & F2 \\ \hline
\endhead
T13 & XGBoost & 0.8869 & 0.3107 & 0.1965 & 0.7602 & 0.2037 \\ \hline
T13 & 固定\textsf{\seqsplit{|Z|$\geq$3}} & 0.0417 & 0.9673 & 0.0672 & 0.4827 & 0.0403 \\ \hline
T18 & XGBoost & 0.8674 & 0.5184 & 0.4246 & 0.8193 & 0.4047 \\ \hline
T18 & 固定\textsf{\seqsplit{|Z|$\geq$3}} & 0.0000 & 0.9409 & 0.0883 & 0.5172 & 0.0000 \\ \hline
T21 & XGBoost & 0.8542 & 0.3022 & 0.0444 & 0.5178 & 0.1171 \\ \hline
T21 & 固定\textsf{\seqsplit{|Z|$\geq$3}} & 0.0000 & 0.9966 & 0.0473 & 0.4851 & 0.0000 \\ \hline
\end{longtable}
\endgroup

XGBoost明显提高T13和T18的灵敏度及PR-AUC，但牺牲特异度和阳性预测值。T21虽然达到较高灵敏度，其PR-AUC仅0.0444、ROC-AUC约0.518，排序能力仍然有限，不能表述为稳定识别。

任一异常概率并集的描述性PR-AUC为0.4458，直接任一异常分类器为0.4264；四个外层折均依据各自训练OOF结果选择概率并集。

\subsection{9. 最终结果}
任一异常概率层的内部外层结果为：灵敏度0.8337、特异度0.4093、阳性预测值0.1525、F2为0.4349、ROC-AUC为0.7756、PR-AUC为0.4622、Brier分数为0.0831。最终决策层将605条记录聚合为147名孕妇，并执行患者级三分流。

\begingroup
\footnotesize
\setlength{\tabcolsep}{1.5pt}
\begin{longtable}{|>{\raggedright\arraybackslash}p{0.280\textwidth}|>{\raggedright\arraybackslash}p{0.280\textwidth}|>{\raggedright\arraybackslash}p{0.280\textwidth}|}
\hline
最终分流 & 患者数 & 含义 \\ \hline
\endfirsthead
\hline
最终分流 & 患者数 & 含义 \\ \hline
\endhead
直接阳性 & 40 & 达到患者级阳性阈值且重复检测满足一致性要求 \\ \hline
建议复测 & 71 & 概率处于复测区、质量异常或仅有T21单独阳性信号 \\ \hline
筛查阴性 & 36 & 未达到直接阳性或复测条件 \\ \hline
\end{longtable}
\endgroup

患者级直接阳性灵敏度为0.500、特异度为0.825、PPV为0.550；直接阳性与复测合并后的覆盖灵敏度为0.909，复测率为0.483，筛查阴性中有4名孕妇曾出现附件异常标签。

2000次患者Bootstrap中，任一异常灵敏度95\%区间约为0.737—0.926、特异度约为0.355—0.462、PR-AUC约为0.311—0.569，说明当前样本下仍存在较大不确定性。

\subsection{10. 补充}
\begin{enumerate}
\item 筛查阳性不等于临床确诊，应结合遗传咨询、羊水穿刺、核型分析或其他诊断性检查。
\item AB列是当前需要复现的附件筛查标签；AE列缺乏真实阴性/阳性变化，不能验证真实疾病诊断能力。
\item 患者级三分流提高了直接阳性的特异度和PPV，但需要48.3\%的复测率，并仍存在4名附件异常标签未覆盖者。
\item T21只有13条阳性记录，Bootstrap区间宽，当前证据不足以宣称稳定预测。
\item 患者级三分流采用复测率不超过55\%的开发约束，外层汇总复测率为48.3\%；该负担必须与覆盖灵敏度0.909同时报告。
\item 当前只有内部嵌套患者级验证，尚无独立医院、真实核型或出生结局外部验证。
\item 最终方法仅用于数学建模和研究性筛查分流，不构成临床诊断或处置意见。
\end{enumerate}

\section{七、补充验证与最终诊断}
\subsection{1. 全部可行组数的嵌套验证}
每个外折将测试孕妇完全留出，在开发患者内部进行3折OOF AFT训练。候选组数取\textsf{\seqsplit{1..floor(n开发/30)}}；约200名开发患者时搜索1—6组，全样本理论上最多8组。各外折的提升轮数仅由对应开发集内层早停确定，依次为\textsf{\seqsplit{150/64/853/113}}。

BMI按数值稳定排序，只允许在相邻不同BMI之间切分。每个区间选择单侧90\%置信下限达到0.95的最早整数孕周，通过动态规划最小化患者加权平均推荐孕周；在最优值0.25周以内选择组数更少的结构。

\begingroup
\footnotesize
\setlength{\tabcolsep}{1.5pt}
\begin{longtable}{|>{\raggedright\arraybackslash}p{0.420\textwidth}|>{\raggedright\arraybackslash}p{0.420\textwidth}|}
\hline
指标 & 结果 \\ \hline
\endfirsthead
\hline
指标 & 结果 \\ \hline
\endhead
选择1组 & 2/4折 \\ \hline
选择2组 & 2/4折 \\ \hline
选择3组及以上 & 0/4折 \\ \hline
平均测试推荐孕周 & 18.949周 \\ \hline
平均测试达标概率 & 0.9556 \\ \hline
均值安全折数 & 4/4 \\ \hline
保守LCB安全折数 & 3/4 \\ \hline
\end{longtable}
\endgroup

最终采用两组可执行政策：\textsf{\seqsplit{BMI<34.357}}第18周、\textsf{\seqsplit{BMI$\geq$34.357}}第22周。固定切点网格中，29.75—34.50为均值与LCB均4/4安全的平台，算法切点34.357位于平台内。

\subsection{2. 患者级筛查三分流}
\begin{itemize}
\item 605条女胎记录聚合为147名孕妇；43名孕妇的重复记录标签不一致，因此患者目标严格定义为“是否曾出现附件筛查异常标签”，不解释为真实疾病。
\item 每个外折只用其余三个外折的OOF患者选择概率阳性阈值和复测下阈值；四折均选择\textsf{\seqsplit{直接阳性阈值0.30/复测下阈值0.15}}。
\item 患者级结果为直接阳性40人、复测71人、阴性36人；直接阳性特异度0.825、PPV 0.550，三分流覆盖灵敏度0.909，复测率0.483，阴性中仍有4名曾出现附件异常标签者。
\item T21单独阳性信号全部进入复测，不直接判为稳定阳性。该方案用于附件筛查标签的研究性分流，不构成临床诊断。
\end{itemize}

\subsection{3. 推荐时点策略敏感性}
\begin{itemize}
\item 5151个准确性—等待—不稳定风险权重点中，15/18、16/18、18/22的胜出比例分别为27.9\%、34.5\%、37.6\%；该比例是离散网格面积近似，不是先验概率。
\item 125套稳定门槛组合中，三套方案的可行率分别为0\%、12.8\%和100\%，说明18/22对稳定门槛最稳健。
\item 严格按每组至少30人且完成4个外折的固定切点表，4/4均值与LCB安全平台为29.75—34.50；算法切点34.357位于平台内。
\end{itemize}

\subsection{4. 主混合模型诊断}
\begin{itemize}
\item 主模型随机截距方差0.1884，未发生随机效应退化；随机截距标准差约0.4153。
\item \textsf{\seqsplit{|标准化残差|>3}}的记录比例为1.386\%，残差与孕周相关约为0，残差与拟合值相关为0.160，提示仍存在轻度结构性误差。
\item 200次患者级Bootstrap全部收敛。孕周线性项区间\textsf{\seqsplit{[0.1357,0.1945]}}、二次项\textsf{\seqsplit{[0.0236,0.0608]}}；BMI项\textsf{\seqsplit{[-0.1292,-0.00054]}}；孕周×BMI交互\textsf{\seqsplit{[-0.02195,0.02899]}}。
\item 结果继续支持孕周正向、BMI负向关联；交互项区间跨0，不支持稳定交互效应。
\end{itemize}

\end{document}
